\section{Solver Methods}

\begin{frame}{Overview: One Engine, Four Methods}
  \textbf{Goal:} minimise $k=\max_e \mathrm{cr}(e)$ (crossings on the worst edge), tie-break on total crossings $\chi$.
  \medskip
  \begin{itemize}
    \item Core engine: two-phase \textbf{Simulated Annealing} after Bianchetti \& Moalic (2025)
      \begin{itemize}
        \item \textbf{Phase 1} minimises total crossings $\chi$ (untangles the drawing)
        \item \textbf{Phase 2} minimises the bottleneck $k$ (dual fitness)
      \end{itemize}
    \item All methods share the same incremental machinery: per-move crossing deltas,
          spatial grids for candidate lookup, strict validity (no vertex on a foreign edge).
  \end{itemize}
  \medskip
  \begin{center}
    \begin{tabular}{llll}
      \toprule
      \textbf{Method} & \textbf{Search} & \textbf{Escapes local optima by} & \textbf{Binary} \\
      \midrule
      \texttt{sa}     & two-phase SA                  & temperature waves          & \texttt{sakgd} \\
      \texttt{staged} & LNS (30\%) $\to$ SA (70\%)    & SA stage after greedy LNS  & \texttt{approach1} \\
      \texttt{ils}    & SA rounds + random kicks      & perturbation between rounds & \texttt{approach1} \\
      \texttt{staged-adaptive} & adaptive-LNS $\to$ SA & growing neighbourhoods     & \texttt{approach1} \\
      \bottomrule
    \end{tabular}
  \end{center}
\end{frame}

\begin{frame}{Method 1: Two-Phase Simulated Annealing (\texttt{sa})}
  \begin{columns}[T]
    \begin{column}{0.55\textwidth}
      \textbf{Move:} pick a vertex (crossing-weighted), propose a new position
      (Gaussian around current, radius shrinks with temperature; 5\% global jumps in phase 1).
      \medskip

      \textbf{Acceptance (Metropolis):} accept if $\Delta E\le 0$, else with prob.\ $e^{-\Delta E/T}$.
      \medskip

      \textbf{Dual fitness in phase 2:}
      \[
        \Delta E =
        \begin{cases}
          \Delta K_{\text{local}} & \text{if } \Delta K_{\text{local}}\neq 0\\[2pt]
          \Delta\chi / \chi & \text{otherwise (tiny tie-break)}
        \end{cases}
      \]
    \end{column}
    \begin{column}{0.42\textwidth}
      \textbf{Temperature waves:}
      \begin{itemize}
        \item inner loop: $T \mathrel{*}= 0.999$ per move until $T_{\lim}$
        \item outer loop: wave restarts from the \emph{best-so-far} solution
              with a 1\% lower starting temperature
      \end{itemize}
      \medskip
      Every accepted improvement of $(k,\chi)$ is snapshotted; the final answer
      is always the best valid layout seen.
    \end{column}
  \end{columns}
\end{frame}

\begin{frame}{Why SA Accepts Worse Moves}
  \begin{columns}[T]
    \begin{column}{0.62\textwidth}
      \resizebox{\linewidth}{!}{%
      \begin{tikzpicture}[
          ball/.style={circle,fill=warm,inner sep=1.7pt},
          gball/.style={circle,fill=accent,inner sep=2.2pt}]
        \draw[muted,line width=1pt] plot[smooth,tension=0.7] coordinates
          {(0,3.0) (0.7,1.8) (1.5,1.2) (2.4,2.6) (3.3,1.0) (4.0,0.4) (5.0,1.3) (6,2.1)};
        \draw[warm,dashed,line width=1pt,-{Stealth}]
          (1.5,1.25) .. controls (2.1,3.3) and (3.0,1.4) .. (3.9,0.55);
        \node[ball]  at (0.3,2.8){};
        \node[ball]  at (1.5,1.2){};
        \node[gball] at (4.0,0.4){};
        \node[font=\scriptsize,align=center,anchor=south]             at (0.55,3.05){start\\(high energy)};
        \node[font=\scriptsize,align=center,text=warm,anchor=south]   at (2.75,3.2){hot: leap over\\the barrier};
        \node[font=\scriptsize,align=center,text=muted,anchor=north]  at (1.5,1.0){local optimum\\(trap)};
        \node[font=\scriptsize,align=center,text=accent,anchor=north] at (4.0,0.2){cold: global\\optimum};
      \end{tikzpicture}}
    \end{column}
    \begin{column}{0.36\textwidth}
      Under pure greedy descent the ball dies in the \textbf{local optimum}.
      \smallskip

      While \emph{hot}, SA tolerates the uphill leap over the barrier; while
      \emph{cold}, it settles and locks in.
      \smallskip

      Accepting $\Delta E>0$ with prob.\ $e^{-\Delta E/T}$ is exactly what
      escapes the trap.
    \end{column}
  \end{columns}
\end{frame}

\begin{frame}{Methods 2--4: LNS, Staged, ILS (\texttt{approach1})}
  \textbf{LNS (destroy \& repair).}
  \begin{itemize}
    \item \emph{Destroy:} BFS-grow a connected neighbourhood from a crossing-weighted seed vertex.
    \item \emph{Repair:} for each vertex, try $\sim$50 candidate positions, commit only the best
          strictly-improving one $\Rightarrow$ effectively greedy hill-climbing: fast descent, but stalls.
  \end{itemize}
  \medskip
  \textbf{Staged (\texttt{lns} $\to$ \texttt{sa}).} Greedy LNS burns down $\chi$ cheaply in the first
  30\% of the budget; SA then escapes the local optimum it left behind. Consistently among the best
  performers on dense graphs.
  \medskip

  \textbf{ILS.} Alternate full SA rounds with random \emph{kicks} (relocate $n/10$ random vertices),
  always restarting from the global best --- a different escape mechanism than temperature.
\end{frame}


\begin{frame}{LNS: Destroy \& Repair}
  \begin{itemize}
    \item \emph{Destroy:} BFS-grow a connected neighbourhood from a crossing-weighted seed vertex.
    \item \emph{Repair:} for each vertex, try $\sim$50 candidate positions, commit only the best
          strictly-improving one $\Rightarrow$ greedy hill-climbing: fast descent, but \textbf{stalls}.
  \end{itemize}
  \vfill
  \centering
  \resizebox{0.92\linewidth}{!}{%
  \begin{tikzpicture}[node distance=10mm,
      b/.style={rectangle,rounded corners,draw=accent,fill=accent!10,align=center,
                font=\scriptsize,inner sep=4pt,text width=3.1cm,minimum height=1.3cm}]
    \node[b](d){\textbf{Destroy}\\BFS-grow a connected neighbourhood from a crossing-weighted seed};
    \node[b,right=of d](r){\textbf{Repair}\\try $\sim$50 spots / vertex; keep the best \emph{strictly-improving} one};
    \node[b,right=of r](s){\textbf{Accept}\\improved $\Rightarrow$ save best; periodically restore best};
    \draw[-{Stealth}](d)--(r);
    \draw[-{Stealth}](r)--(s);
    \draw[-{Stealth}](s.south) .. controls +(0,-1.1) and +(0,-1.1) .. (d.south);
    \node[font=\scriptsize,text=warm]  at ($(r.north)+(0,0.45)$){no uphill moves $\Rightarrow$ stalls on plateaus};
    \node[font=\scriptsize,text=muted] at ($(r.south)+(0,-1.15)$){repeat until $k=0$ or out of time};
  \end{tikzpicture}}
\end{frame}

\begin{frame}{Staged \& ILS}
  \textbf{Staged} (\texttt{lns}$\to$\texttt{sa}): greedy LNS burns down $\chi$ in the first 30\% of the
  budget; SA then escapes the local optimum it left behind. \emph{Consistently among the best on dense graphs.}
  \vfill
  \begin{columns}[T]
    \begin{column}{0.66\textwidth}
      \begin{tikzpicture}
        \begin{axis}[width=\linewidth,height=5cm,
            xlabel={budget used (\%)},ylabel={total crossings $\chi$},
            xmin=0,xmax=100,ymin=0,ymax=108,legend pos=north east,
            tick label style={font=\scriptsize},label style={font=\scriptsize}]
          \addplot[teal,very thick] coordinates
            {(0,100)(3,70)(6,52)(10,40)(15,33)(20,30)(25,28.5)(30,28)};
          \addlegendentry{LNS (first 30\%)}
          \addplot[warm,very thick] coordinates
            {(30,28)(35,26)(45,20)(55,15)(70,11)(85,9)(100,8)};
          \addlegendentry{SA (rest)}
          \draw[muted,dashed] (axis cs:30,0)--(axis cs:30,108);
          \node[font=\scriptsize,text=teal,anchor=west]            at (axis cs:6,64){fast greedy descent};
          \node[font=\scriptsize,text=muted,anchor=south]          at (axis cs:23,31){stalls};
          \node[font=\scriptsize,text=warm,anchor=west,align=left] at (axis cs:48,33){SA escapes the\\plateau};
          \node[font=\scriptsize,text=muted,anchor=north]          at (axis cs:30,107){handoff};
        \end{axis}
      \end{tikzpicture}
    \end{column}
    \begin{column}{0.32\textwidth}
      \footnotesize
      \textbf{ILS:} alternate full SA rounds with random \emph{kicks} (relocate $n/10$ vertices),
      always restarting from the global best --- a different escape mechanism than temperature.
    \end{column}
  \end{columns}
\end{frame}